\section{The critical field rises convexly with the coupling}
\label{sec:3}

\subsection{The onset is a Sturm--Liouville problem on the brane}
\label{sec:3.1}

In the probe limit the instability is the one of Nielsen and Olesen: the
charged combination $W = A^1 + iA^2$ of the gauge field is a vector boson of
gyromagnetic ratio $g = 2$, and in the lowest Landau level its
magnetic-moment coupling to the background field is strong enough to make it
tachyonic. As shown in~\cite{Ammon:2011je}, the mode that condenses first is
the polarised lowest Landau level
\begin{equation}
W_x = w(r)\,\psi_0(x,y)~, \qquad W_y = -i\,W_x~,
\end{equation}
where $\psi_0$ is any function in the lowest Landau level of the charge-one
field $A^3 = Bx\,dy$, and $w(r)$ is the radial profile. In the boundary
theory what condenses is the charged spatial current
$\langle J^1_i \pm iJ^2_i\rangle$, modulated in the plane. The same
instability for a charged vector of general mass and gyromagnetic ratio,
which reduces to the $SU(2)$ case at zero mass and $g = 2$, was studied in
the probe limit in~\cite{Cai:2013pda,Cai:2013kaa}. In the gauge $W_r = 0$,
$W_t$ and $W_z$ vanish identically for this mode, and on the magnetic brane
the same ansatz reduces the linearised Yang--Mills equations exactly to a
Sturm--Liouville problem in $r$ with the field as the eigenvalue,
\begin{equation}
\big(P\,w'\big)' + Q\,w = 0~, \qquad P = U\,e^{Z}~, \qquad Q = B\,e^{Z-2V}~.
\end{equation}
One feature of this equation matters for what follows. On the polarised lowest Landau level
the transverse kinetic term of $W$ vanishes identically, since
$F^W_{xy} = 0$ for this polarisation, so the only term without a radial
derivative is the non-abelian magnetic-moment coupling. This is the $g = 2$
mechanism, and it carries a single factor $e^{-2V}$: the moment term samples
the warping of the flux plane and nothing else, which is what makes the
field drop out in the throat (section~\ref{sec:4.2}). The onset is
determined by this one equation on the full brane (the truncation is
consistent: at linear order in $w$ nothing sources the metric, the abelian
field $A^3$ or the other colour components;
appendix~\ref{app:charged-zero-mode}). In the probe limit $\beta \to 0$ the equation
reduces to the probe problem, $P \to f/u$ and $Q \to B/u$ in the coordinate
$u = 1/r$, with $f = 1 - u^4$ the blackening factor of
$\mathrm{AdS}_5$--Schwarzschild, and the same critical field $B_c u_H^2 = 5.13126764$.

The boundary conditions are those of the probe calculation, regularity at
the horizon and no source at the boundary, and the brane's logarithmic
asymptotics do not obstruct them (appendix~\ref{app:charged-zero-mode}).
Since $\psi_0$ is
arbitrary within the lowest Landau level, the onset is infinitely
degenerate, and the lattice is selected only at quartic order, which is the
subject of section~\ref{sec:5}.

\subsection{The onset curve is a single sweep in \texorpdfstring{$\beta$}{beta}}
\label{sec:3.2}

At fixed $\alpha$ the critical field is the smallest eigenvalue $B$ of the
Sturm--Liouville problem on the brane at the same $B$. This looks like a
self-consistency problem, but it needs no iteration: the background depends
on $\alpha$ and $B$ only through $\beta$, and at fixed $\beta$, $B$ enters
the operator only as the eigenvalue. We therefore construct the brane at a given
$\beta$, solve for its lowest
eigenvalue $B_{\rm eig}(\beta)$, and read off the coupling at which this
brane is the self-consistent onset,
\begin{equation}
B_c = B_{\rm eig}(\beta)~, \qquad \alpha(\beta) = \frac{\beta}{B_{\rm eig}(\beta)^2}~.
\end{equation}

\subsection{Backreaction stabilises the normal phase without saturating}
\label{sec:3.3}

The critical field at fixed temperature is shown in figure~\ref{fig:1}, in
units of the horizon radius, that is $B_c/T^2 = \pi^2 B_c u_H^2$. We solved
the sweep on a dense grid in $\beta$ up to $\beta = 45$, where
$\alpha = 0.1745$ and $B_c u_H^2 = 16.06$, and on a sparse sequence of 23
values from $\beta = 1$ to $10^{11}$, which we call the ladder and its
members rungs; the rungs beyond $\beta = 45$ are used for the approach
to $\alpha_\star$ in section~\ref{sec:4.6}. The curve is
monotone and convex, and its tangent at $\alpha = 0$ is already ten percent
off at $\alpha \approx 0.08$. Backreaction thus stabilises the normal phase: at a
given temperature the magnetised plasma needs a larger field to condense
when the field is allowed to deform the geometry, and the effect grows
faster than linearly. That backreaction works against the instability, and
removes it at zero temperature beyond a critical coupling, was found
in~\cite{Wong:2013rda} in four bulk dimensions; the exact curve shows how
$B_c$ diverges on the approach to $\alpha_\star$. Along the ladder $B_c$ rises by more than three orders
of magnitude while $\alpha$ moves from $0.4$ to $0.59$,
which is the first sign that the curve has a vertical asymptote.
